Standard question-answering benchmarks have driven strong progress by asking whether a system can produce a correct answer for a given question or context, as in Natural Questions and SQuAD-style reading comprehension \citep{kwiatkowski-etal-2019-natural,rajpurkar-etal-2018-know}. Broader multitask and holistic evaluation suites similarly foreground final-answer correctness across domains and scenarios \citep{hendrycks-etal-2021-mmlu,srivastava-etal-2023-bigbench,liang-etal-2023-helm}. Our setting starts from a different failure boundary: before judging the answer text, the assistant must decide whether answering is the right next action at all.

Several lines of work study cases where the user query is underspecified or ambiguous. Clarification work asks systems to elicit missing user intent in information-seeking conversations \citep{aliannejadi-etal-2019-asking}; AmbigQA and ASQA require systems to resolve or cover multiple plausible interpretations of an ambiguous question \citep{min-etal-2020-ambigqa,stelmakh-etal-2022-asqa}. These tasks motivate our \texttt{ask} action, but they do not by themselves cover cases where the user premise is false, stale, or contradicted by retrieved evidence.

False-premise, fact-verification, and truthfulness benchmarks target a complementary defect. CREPE studies open-domain questions containing false presuppositions and asks models to identify and correct them \citep{yu-etal-2023-crepe}; FEVER evaluates whether claims are supported or refuted by evidence \citep{thorne-etal-2018-fever}; TruthfulQA probes whether models reproduce common false beliefs \citep{lin-etal-2022-truthfulqa}. Temporal and situated QA work further shows that models and retrieval-augmented systems can return outdated or context-insensitive answers when the world, location, or time changes \citep{kasai-etal-2023-realtimeqa,vu-etal-2023-freshllms,zhang-choi-2021-situatedqa}. We treat false and stale premises as action-selection cases: the useful response often begins by challenging the premise rather than answering through it.

Other work studies uncertainty, abstention, and conflict in the evidence. SQuAD 2.0 introduced adversarially written unanswerable questions for reading comprehension \citep{rajpurkar-etal-2018-know}. Selective QA under domain shift evaluates whether QA systems can abstain on likely errors while still answering as many questions as possible \citep{kamath-etal-2020-selective}, connecting to broader selective-classification and calibration work \citep{geifman-el-yaniv-2017-selective,guo-etal-2017-calibration}. Recent retrieval-conflict datasets and adaptive QA settings examine how systems answer when retrieved contexts disagree or require source-aware reasoning \citep{shaier-etal-2024-adaptive,liu-etal-2025-open}. This conflict setting is downstream of the broader retrieval-augmented and evidence-grounded QA line \citep{guu-etal-2020-realm,lewis-etal-2020-rag,izacard-grave-2021-fid,nakano-etal-2021-webgpt,menick-etal-2022-quotes}. Work on language-model self-evaluation, calibration, and uncertainty asks whether models can predict when their own answers are likely to be correct \citep{kadavath-etal-2022-language,manakul-etal-2023-selfcheckgpt,kuhn-etal-2023-semantic}. Our benchmark differs by mapping these conditions into the same action space and scoring the chosen next action directly: \texttt{abstain} is one option, but false or stale premises often require \texttt{challenge}, and ambiguous intent requires \texttt{ask}.

Taken together, prior work establishes important specialized tasks: clarify ambiguity, correct false assumptions, answer time-sensitive questions, abstain under missing support, and handle conflicting evidence. This paper keeps the benchmark contribution focused: evaluating those behaviors through one assistant policy problem, so a model must choose among \texttt{answer}, \texttt{ask}, \texttt{challenge}, and \texttt{abstain} before it can safely generate the final text. That boundary is the difference between shared ontology and mixed-task scoring.
